% C2 - mTLS Service Mesh
% Documento LaTeX independiente para incluir con % C2 - mTLS Service Mesh
% Documento LaTeX independiente para incluir con % C2 - mTLS Service Mesh
% Documento LaTeX independiente para incluir con % C2 - mTLS Service Mesh
% Documento LaTeX independiente para incluir con \input{documentecionFinal/C2_mtls_service_mesh.tex}

\section{Control C2: mTLS Service Mesh}
\label{sec:c2-mtls-service-mesh}

El control C2 eval\'ua el impacto de habilitar mTLS entre microservicios con dos mallas de servicio, comparadas contra un baseline sin sidecars.

\subsection{Configuraci\'on aplicada}

\begin{itemize}
  \item \textbf{Baseline:} se usa el despliegue est\'andar de servicios (sin malla), con comunicaci\'on interna directa entre \texttt{api-service}, \texttt{data-service} y \texttt{auth-service}.

  \item \textbf{Istio mTLS:} se aplica el manifiesto con inyecci\'on de sidecar de Istio habilitada (namespace y pods con inyecci\'on activa), manteniendo los mismos endpoints funcionales de los servicios.

  \item \textbf{Linkerd mTLS:} se aplica el manifiesto con anotaci\'on \texttt{linkerd.io/inject: enabled} en los deployments de los microservicios para habilitar el proxy sidecar.
\end{itemize}

En t\'erminos pr\'acticos, C2 mantiene la misma topolog\'ia funcional y cambia solo la capa de transporte/sidecar para medir el costo operacional de mTLS. En S2, la entrada de prueba puede provenir de ingress o port-forward; la variaci\'on experimental de C2 est\'a en la capa inter-servicio.


\section{Control C2: mTLS Service Mesh}
\label{sec:c2-mtls-service-mesh}

El control C2 eval\'ua el impacto de habilitar mTLS entre microservicios con dos mallas de servicio, comparadas contra un baseline sin sidecars.

\subsection{Configuraci\'on aplicada}

\begin{itemize}
  \item \textbf{Baseline:} se usa el despliegue est\'andar de servicios (sin malla), con comunicaci\'on interna directa entre \texttt{api-service}, \texttt{data-service} y \texttt{auth-service}.

  \item \textbf{Istio mTLS:} se aplica el manifiesto con inyecci\'on de sidecar de Istio habilitada (namespace y pods con inyecci\'on activa), manteniendo los mismos endpoints funcionales de los servicios.

  \item \textbf{Linkerd mTLS:} se aplica el manifiesto con anotaci\'on \texttt{linkerd.io/inject: enabled} en los deployments de los microservicios para habilitar el proxy sidecar.
\end{itemize}

En t\'erminos pr\'acticos, C2 mantiene la misma topolog\'ia funcional y cambia solo la capa de transporte/sidecar para medir el costo operacional de mTLS. En S2, la entrada de prueba puede provenir de ingress o port-forward; la variaci\'on experimental de C2 est\'a en la capa inter-servicio.


\section{Control C2: mTLS Service Mesh}
\label{sec:c2-mtls-service-mesh}

El control C2 eval\'ua el impacto de habilitar mTLS entre microservicios con dos mallas de servicio, comparadas contra un baseline sin sidecars.

\subsection{Configuraci\'on aplicada}

\begin{itemize}
  \item \textbf{Baseline:} se usa el despliegue est\'andar de servicios (sin malla), con comunicaci\'on interna directa entre \texttt{api-service}, \texttt{data-service} y \texttt{auth-service}.

  \item \textbf{Istio mTLS:} se aplica el manifiesto con inyecci\'on de sidecar de Istio habilitada (namespace y pods con inyecci\'on activa), manteniendo los mismos endpoints funcionales de los servicios.

  \item \textbf{Linkerd mTLS:} se aplica el manifiesto con anotaci\'on \texttt{linkerd.io/inject: enabled} en los deployments de los microservicios para habilitar el proxy sidecar.
\end{itemize}

En t\'erminos pr\'acticos, C2 mantiene la misma topolog\'ia funcional y cambia solo la capa de transporte/sidecar para medir el costo operacional de mTLS. En S2, la entrada de prueba puede provenir de ingress o port-forward; la variaci\'on experimental de C2 est\'a en la capa inter-servicio.


\section{Control C2: mTLS Service Mesh}
\label{sec:c2-mtls-service-mesh}

El control C2 eval\'ua el impacto de habilitar mTLS entre microservicios con dos mallas de servicio, comparadas contra un baseline sin sidecars.

\subsection{Configuraci\'on aplicada}

\begin{itemize}
  \item \textbf{Baseline:} se usa el despliegue est\'andar de servicios (sin malla), con comunicaci\'on interna directa entre \texttt{api-service}, \texttt{data-service} y \texttt{auth-service}.

  \item \textbf{Istio mTLS:} se aplica el manifiesto con inyecci\'on de sidecar de Istio habilitada (namespace y pods con inyecci\'on activa), manteniendo los mismos endpoints funcionales de los servicios.

  \item \textbf{Linkerd mTLS:} se aplica el manifiesto con anotaci\'on \texttt{linkerd.io/inject: enabled} en los deployments de los microservicios para habilitar el proxy sidecar.
\end{itemize}

En t\'erminos pr\'acticos, C2 mantiene la misma topolog\'ia funcional y cambia solo la capa de transporte/sidecar para medir el costo operacional de mTLS. En S2, la entrada de prueba puede provenir de ingress o port-forward; la variaci\'on experimental de C2 est\'a en la capa inter-servicio.
